% !PRINTED 322
\weberpage{322}

A fundamental discriminant is divisible by no odd square number. If, however, $\Delta$ is divisible by 4, then
\weq{3}{\Delta \equiv 8,\ 12 \quad (\mathrm{mod}\ 16)}
must hold. For were $\Delta \equiv 0, 4$ $(\mathrm{mod}\ 16)$, then after cancellation of the factor 4 a discriminant would remain.

In order, then, to find the associated fundamental discriminant for a given $D$, one has first to separate off the greatest odd square number, and then in addition so high a power of 4 that $\Delta$ becomes either odd and $\equiv 1$ (mod 4), or $\equiv 8$, or $\equiv 12$ (mod 16).

4. \gesperrt{A proper discriminant which is not decomposable into factors that are themselves again discriminants is called a prime discriminant.}

If $p$ is a natural odd prime, and if the sign $\pm$ is so determined that $\pm p \equiv 1$ (mod 4), then the prime discriminants are the following:
\weq{4}{\pm p,\quad -4,\quad +8,\quad -8.}

5. \gesperrt{Every fundamental discriminant can be decomposed in one and only one way into prime discriminants.}

For if one first separates off from $\Delta$ all the factors $\pm p$, there remains only one of the numbers 1, $-4$, $+8$, $-8$.

\begin{center}
\textbf{§ 85. The extended \gesperrt{Legendre-Jacobi} symbol.}
\end{center}

\gesperrt{The Legendre symbol} $\left(\dfrac{m}{p}\right)$ has, when $m$ is an arbitrary positive or negative number different from zero and $p$ is a natural prime not dividing $m$, the value $+1$ if $m$ is a quadratic residue, and the value $-1$ if $m$ is a quadratic non-residue of $p$ (Bd. I, § 145).

The meaning of the symbol has been extended by \gesperrt{Jacobi} in such a way that a composite number may also be put for $p$. If, however, $m$ is a discriminant, then a still further reaching generalization is sometimes to be recommended, and this we must now set forth\footnote{\gesperrt{Kronecker}, Berliner Sitzungsberichte, 30. Juli 1885. \gesperrt{Dirichlet-Dedekind}, Vorlesungen über Zahlentheorie, 4. Aufl., § 186. H.\ \gesperrt{Weber}, Göttinger Nachrichten, Januar 1893.}.
